% ============================================================
% Results — the application to the real sample and what was
% found. Report measurements; save interpretation for the
% Discussion. Large measurement tables go in the appendix; a
% summary table stays here.
% ============================================================
\section{Results}
\label{sec:results}

\subsection{Measurements}
\label{subsec:measurements}

% The primary outputs: a summary table of measured quantities
% (full catalog in the appendix or a machine-readable table —
% AAS journals host data behind the article):
% \begin{table}[htbp]
%     \caption{Summary of measurements}
%     \label{tab:summary}
%     \centering
%     \begin{tabular}{lcccc}
%         \hline\hline
%         Sample & $N$ & Quantity A & Quantity B & Units \\
%         \hline
%              ... & ... & ... & ... & ... \\
%         \hline
%     \end{tabular}
% \end{table}

\subsection{Key Findings}
\label{subsec:findings}

% One subsection per finding. Each: the correlation/detection,
% its statistical significance (with the actual numbers and
% uncertainties), and the figure that shows it. Scaling-relations
% and their scatter are a common currency — state outliers, not
% just the trend.

% State systematic uncertainty alongside statistical where they
% compete ("a 0.05 dex systematic from X, comparable to the
% 0.04 dex statistical scatter").
